% !TEX root = ../main.tex

この部では、圏論の中心概念を、できるだけソフトウェアに近い形で導入する。圏論という名前からは、非常に抽象的な数学を想像するかもしれない。しかし本書では、最初から一般の圏を厳密に分類するのではなく、まず「型と関数の世界」を使って考える。型を対象、関数を射、関数合成を射の合成として見ると、圏論の多くの概念は、日々のプログラミングで行っていることの整理として現れる。

最初に扱うのは、圏そのものである。圏とは、合成できるものの集まりである。何もしない処理である恒等射があり、処理をつなぐ合成があり、その合成が自然な法則に従う。この考え方は、関数パイプライン、ETL、middleware、adapter chain など、さまざまな設計に現れる。

次に、同型を扱う。同型とは、情報を失わずに相互変換できる関係である。ソフトウェアでは、スキーマ移行、シリアライズとデシリアライズ、DTO変換、旧APIと新APIの互換性などで重要になる。ただし、\texttt{A → B} と \texttt{B → A} という二つの関数があるだけでは同型ではない。戻したときに本当に元に戻ること、つまり round-trip 性質が必要である。この性質をLeanで証明することが、後半のマイグレーション章の基礎になる。

さらに、積と余積、関手、自然変換、モノイド、モナドを導入する。積はレコードやペア、余積は \texttt{Either} やエラー分岐、関手は \texttt{List.map} や \texttt{Option.map}、自然変換は型に依存しない adapter、モノイドはログや設定の結合、モナドは失敗や状態を含む処理の合成として読むことができる。

この部では、各概念を説明した直後にLeanで最小モデルを示す。Mathlibの高度な圏論ライブラリにはすぐ進まず、まず自分で小さく定義する。抽象化の威力を理解するためには、最初に自分の手で「型と関数の圏」を触ることが重要だからである。

\section*{この部で扱うこと}

\begin{itemize}
\item 対象、射、合成、恒等射
\item 同型と round-trip law
\item 積と余積
\item 関手と \texttt{map}
\item 自然変換と自然性の四角形
\item モノイドと結合律
\item モナドと effectful computation の合成
\item 各概念のLeanによる最小表現
\end{itemize}

\section*{この部で扱わないこと}

\begin{itemize}
\item 一般圏論の高度な定理
\item 米田の補題の完全な証明
\item Kan拡張
\item トポス理論
\item \(\infty\)圏論
\end{itemize}

\section*{次の部への接続}

次の部では、ここで出てきた概念を実際に証明で使うための道具箱を整える。等式推論、場合分け、帰納法、関数外延性、Galois接続、Pullback などを、定理として暗記するのではなく、仕様や移行の安全性を確認するための実用的な道具として学ぶ。
